\section{Divisibility in a weighted companion algebra}
\label{sec:weighted-companion-saturation}

We prove the coefficient divisibility statement used at every finite
stage of the source construction.  Completions and the source division
are included in the statement.

\subsection{The source and its coefficient representation}

Let \(L/\mathbb Q_p\) be finite, with ring of integers \(O\), and put
\[
 R=O[[\gamma-1]][1/p].
\]
Let \(W=R^r\), with its regular \(\Gamma=\langle\gamma\rangle\)-action.
For \(2\leq j\leq d\), let \(V_j\) be a finite-dimensional
\(L\)-space with continuous \(\Gamma\)-action.  A zero space is allowed.
Form the completed free weighted associative algebra
\[
 \mathcal A^{\leq d}
   =\widehat T_L(W,V_2,\ldots,V_d),
 \qquad \operatorname{wt}(W)=1,\quad
 \operatorname{wt}(V_j)=j.
\]
In each fixed degree, the completion is obtained from compact
\(O\)-lattices and then by inverting \(p\).  In particular an entire
coefficient array has one scalar denominator.

Write
\[
 S_j=(\mathcal A^{<j})_j,\qquad
 E_j=(\mathcal A^{\leq j})_j=S_j\oplus V_j
\]
as \(L\)-spaces.  The action on the generators has the form
\begin{equation}\label{eq:wcs-action}
 \gamma v=A_jv+C_j(v),\qquad C_j(v)\in S_j,
\end{equation}
where \(A_j\) is the given finite action.  Thus it defines an actual
weight-preserving algebra automorphism, triangular in generator
weight.  We require that each actual module \(E_j\) is \(R\)-torsion-free.
For the Lie assertion below, the companions \(C_j(v)\) are primitive;
equivalently, \(\gamma\) is a Hopf automorphism for the coproduct that
makes all generators primitive.  These properties will be verified
for the arithmetic source modules when the theorem is applied.

At a fixed degree only finitely many companion coefficients occur.
Choose a free compact lattice in \(W\), and scale the finitely many
generator lattices successively so that the companions for both
\(\gamma\) and \(\gamma^{-1}\) are integral.  This gives a compact
\(\Gamma\)-stable model.  This operation does not assert integral
\(\Lambda\)-freeness of an arithmetic lattice originally giving \(W\).

Set \(z_i=1+X_i\).  The regular coefficient map on \(W\) is the
identity map to the one-variable power-series arrays, with one
coordinate for each of its \(r\) colors.  Products concatenate
ordered variable blocks.  If a weight-\(j\) generator occupies a
contiguous block \(I\), its coefficient is defined by the exact
equation
\begin{equation}\label{eq:wcs-resolvent}
 (z_I I-A_j)\,\iota(v)=\iota(C_j(v)),
 \qquad z_I=\prod_{i\in I}z_i.
\end{equation}
Choose a polynomial \(q_j\) annihilating \(A_j\).  The resulting map in
degree \(n\) is
\[
 \iota_n:(\mathcal A^{\leq d})_n
 \longrightarrow \mathcal B_n^{\,r^n},\qquad
 \mathcal B_n=
 O[[X_1,\ldots,X_n]]
 [1/p,q_j(z_I)^{-1}:|I|=j,\ I\text{ contiguous}].
\]
In this expression a localization uses a common polynomial
denominator for the whole array.  For \(n>d\), all these intervals
are proper.  The action of \(\gamma\) on the target is multiplication
by \(z_1\cdots z_n\).

\begin{theorem}\label{thm:weighted-companion-saturation}
The maps \(\iota_n\) are injective.  For every \(n>d\),
\(\beta\in1+\mathfrak m_L\), and \(k\geq1\), put
\(D_\beta=z_1\cdots z_n-\beta\).  Then
\begin{equation}\label{eq:wcs-saturation}
 \iota_n((\mathcal A^{\leq d})_n)
       \cap D_\beta^k\mathcal B_n^{\,r^n}
 =\iota_n\bigl((\gamma-\beta)^k
                       (\mathcal A^{\leq d})_n\bigr).
\end{equation}
When the companions are primitive, the same equality holds on
the completed free weighted Lie subalgebra.
\end{theorem}

Degree \(n\leq d\) is deliberately absent.  A generator of weight
\(n\) can have its full total divisor inverted in
\eqref{eq:wcs-resolvent}; its source coinvariants need not vanish.

\subsection{Bounded local coordinates and balanced tensors}

\begin{lemma}\label{lem:wcs-local-coordinates}
Each injected fixed-degree source module is locally pro-free
over \(R\), with finitely many coordinate lines rescaled by powers
of a local parameter.  For an extension
\[
 0\longrightarrow S_j\longrightarrow E_j\longrightarrow V_j
       \longrightarrow0
\]
whose old coefficient map is saturated on its total divisor, these
coordinate changes may be chosen simultaneously in the source and
on its actual coefficient columns.  They use finitely many
coefficient functionals and one auxiliary denominator.
\end{lemma}

\begin{proof}
First consider the closed ideal in a fixed-degree source consisting
of words containing at least one \(W\)-letter.  Filter it by the
number of \(W\)-letters.  A companion substitution never decreases
this number.  Its graded pieces are finite sums of a tensor power
of \(W\) and finite modules of zero-\(W\) gap words.

On the first regular \(W\)-factor the exact map
\[
       [g]\otimes m\longmapsto[g]\otimes g^{-1}m
\]
turns the diagonal action into the regular action on that factor.
The map and its inverse are defined on compatible finite
group-algebra quotients and finite invariant coefficient quotients.
Thus each graded piece is compact pro-free over \(O[[\Gamma]]\).

A continuous compact surjection onto such a module has a
continuous module section.  Indeed, its convergent free basis
has lifts tending to zero: choose them through a decreasing
sequence of open submodules, using openness of the compact
surjection.  The series formed from those lifts converges on
every finite quotient and defines the section.  Splitting the
finite filtration proves that the ideal is pro-free.  The
quotient consists of finitely many zero-\(W\) words and is finite
over \(O\).

After enlarging \(L\), fix a root \(\alpha\), and put
\(f=\gamma-\alpha\).  For the top lifts of the Jordan blocks of
the finite quotient write
\[
                    f^{r_i}y_i=C_i
\]
in the pro-free submodule.  The classes of the \(C_i\) modulo
\(f\) are independent over \(L\).  A relation
\(\sum c_iC_i=fs\) would make
\[
 f\left(\sum c_i f^{r_i-1}y_i-s\right)=0.
\]
Torsion-freeness of the full source would force the parenthesized
element to vanish, contradicting its nonzero Jordan-socle image.

Choose finitely many coefficient coordinates giving an invertible
minor on these residues.  Its determinant \(g\) is nonzero at
\(f=0\).  After inverting \(g\), a finite-rank change of coordinates
sends the full \(f\)-dependent columns \(C_i\) to coordinate
vectors.  The local source is obtained by rescaling these finitely
many lines by \(f^{-r_i}\).  Every operation has one common
denominator.

For the stated companion extension, old total saturation injects
\(S_j/fS_j\) into the actual specialized coefficient module.
Multiply the companion columns by a common proper-interval
denominator and choose the same finite minor there.  This gives
bounded source and coefficient projectors
\(\ell_i(C_a)=\delta_{ia}\), as well as their complementary
projector.  Locally,
\begin{equation}\label{eq:wcs-polar-coordinates}
 E_j=S_j^\circ\oplus\bigoplus_i R_f f^{-r_i}C_i,\qquad
 S_j=S_j^\circ\oplus\bigoplus_i R_f C_i .
\end{equation}
Here \(R_f\) indicates the local ring with the one auxiliary
denominator allowed.  The projectors commute with \(f\).  The
regular complement inherits every power of the old saturation.
\end{proof}

An inverse of \(g\) need not have a bounded infinite Taylor
expansion in \(f\).  The proof only uses a fixed finite jet at
each step.  Its coefficients have a common bound, obtained by
inverting a fixed power of \(g(0)\).  This is why the local
coordinate argument remains valid for bounded completions.

\begin{lemma}\label{lem:wcs-balanced}
For two previously constructed degree modules \(E_i,E_j\),
whose smaller alphabets have the total saturation of
Theorem~\ref{thm:weighted-companion-saturation}, the actual
balanced coefficient map is injective:
\[
 (E_i\widehat\otimes E_j)/(\gamma\otimes\gamma-\beta)
   =E_i\widehat\otimes_{R,\tau_\beta}E_j
       \longrightarrow \mathcal B_{i+j}/D_\beta,
 \qquad \tau_\beta(\gamma)=\beta\gamma^{-1}.
\]
All tensor products and source quotients use the compact
models followed by inversion of \(p\).
\end{lemma}

\begin{proof}
The displayed source equality is the quotient by the balancing
relation on pure tensors, followed by its continuous extension.
Choose the total variable \(r\) of the first block; the second
total is then \(\beta/r\).  The two sets of internal variables
remain disjoint.

First take the two old modules \(S_i,S_j\).  At \(r=\alpha\),
their maps modulo \(f=r-\alpha\) are injective by old total
saturation, with the second root transported to \(\beta/\alpha\).
The completed tensor of these injections is injective.
To check this last statement within the bounded convention,
clear the fixed denominators.  A compact torsion-free
\(O\)-module is a product of copies of \(O\); completed tensoring
with such a product gives a product of the original injection.
The target is exactly the bounded coefficient module in the
disjoint internal variables.

In the local coordinates of Lemma~\ref{lem:wcs-local-coordinates},
a nonzero balanced tensor has a least \(f\)-adic order.
Monic division by that power is bounded.  Its nonzero reduction
would have nonzero coefficient by the preceding paragraph.
This proves the old balanced injection, and the same argument
proves its \(f\)-saturation.

Now replace \(S_i,S_j\) by \(E_i,E_j\).  At a root of either
finite companion polynomial, use
\eqref{eq:wcs-polar-coordinates} on both factors.  Their
independent companion columns give finite local coordinate
changes in the source and coefficient modules.  The tensor
coordinates are rescaled by the sums of their two exponents;
unchanged coordinates have exponent zero.  The old injection
therefore remains injective on each rescaled block.

The possible global kernel is killed by the product of the
two finite companion polynomials, the second transported by
\(\tau_\beta\).  Its local vanishing at their finitely many
roots implies global vanishing by Bezout.  These are finite
polynomials and finite coordinate changes; no separate
denominator has been assigned to an infinite coefficient.
\end{proof}

\subsection{Interval poles and the boundary bound}

On \(z_1\cdots z_n=\beta\), two proper contiguous-interval
divisors can coincide only when their intervals and roots
coincide, or when their intervals are complementary prefix
and suffix intervals and their roots multiply to \(\beta\).
Indeed, after fixing the full product and one interval product,
a second monomial is constant exactly when its indicator vector
lies in the span of those two indicator vectors.  For zero-one
indicators this gives the interval, its complement, the empty
set, or the full set.  The complement is contiguous precisely
for a prefix or suffix.

Recursively expanding \eqref{eq:wcs-resolvent} produces finitely
many companion-tree patterns at a fixed weighted degree.
Each finite vertex occupies a contiguous interval and supplies
the denominator of that interval's total.  Ancestor weights
strictly increase.  The leaf coefficients remain completed
bounded \(W\)-arrays; only the finite list of token patterns
is expanded.

\begin{lemma}\label{lem:wcs-boundary-bound}
Put \(n=d+s<2d\), \(1\leq s<d\), and work on the total
hypersurface at
\[
 r=z_1\cdots z_d,\qquad t=z_{d+1}\cdots z_n=\beta/r,
 \qquad f=r-\alpha.
\]
Consider all source terms except those beginning with a new
weight-\(d\) token.  In the local suffix coordinates
\eqref{eq:wcs-polar-coordinates}, their negative principal
parts lie only in the suffix companion columns
\(f^{-q_a}D_a\), and their pole order on that column is at
most \(q_a\).  They have no negative principal part in
the suffix regular complement.
\end{lemma}

\begin{proof}
By the interval calculation, an \(f\)-pole of one of these
terms must come from a vertex occupying the entire complementary
suffix of length \(s\).  No other term has a vertex occupying
the full prefix of length \(d\).  This suffix vertex has
weight \(s\), occurs once in its tree, and has root
\(\beta/\alpha\).

Every ancestor of this vertex is a suffix of length \(j>s\).
On the total hypersurface its total is exactly
\begin{equation}\label{eq:wcs-ancestor}
 z_{n-j+1}\cdots z_n
       =\frac{\beta}{z_1\cdots z_{n-j}}.
\end{equation}
Choose \(z_1,\ldots,z_{d-1}\) as the internal prefix
coordinates and \(r\) as its total.  Since \(n-j<d\),
\eqref{eq:wcs-ancestor} is a function of proper prefix
coordinates, independent of \(f\).  Every coefficient
outside the suffix vertex depends on prefix variables
and the suffix total only.  Its denominators are regular
at \(f\), and its ancestor denominators are the fixed
proper-prefix factors in \eqref{eq:wcs-ancestor}.

The vertex itself supplies an exact suffix column
\((t-\beta/\alpha)^{-q_a}D_a(t)\).  The factor
\(t-\beta/\alpha\) is a unit multiple of \(f\).
Thus the negative principal part is a sum of
\[
 h_a(f,\text{prefix})\,f^{-q_a}D_a(f,\text{suffix}),
 \qquad h_a\text{ regular at }f.
\]
This includes an old token crossing the cut and a new suffix
token whose companion contains the suffix vertex.

There are only finitely many negative orders.  Discard all
\(f\)-regular terms before applying any coefficient
projector.  Their possibly mixed prefix-suffix denominators
therefore never enter the extraction space.  For the remaining
polar terms, take only the finite jets required by their
orders.  A fixed power of each nonzero constant denominator
clears these jets.  The resulting denominators separate into
a common prefix denominator and a common suffix denominator.
The principal-part identity consequently lies in the
completed tensor of the two separate coefficient modules.
The full suffix projectors of
Lemma~\ref{lem:wcs-local-coordinates} apply to that identity
and give the stated bounds.  Using full columns, rather
than just their residues, handles derivatives of unequal
Jordan blocks.
\end{proof}

\subsection{Proof of the saturation theorem}

\begin{proof}[Proof of Theorem~\ref{thm:weighted-companion-saturation}]
We induct on the largest generator weight, retaining the
assertion for every smaller alphabet.  For the regular
alphabet \(W\), change one variable to the full product.
Monic power-series division in that coordinate proves the
statement with a common coefficient bound.

Write \(\mathcal A=\mathcal A^{<d}\),
\(\mathcal B=\mathcal A^{\leq d}\).  Old total saturation
in degree \(d\) and torsion-freeness of \(E_d\) give the
independent full companion residues of Lemma~\ref{lem:wcs-local-coordinates}.  The ordinary
coefficient map of \(\mathcal B\) is injective: take the
largest initial token weight, extract its highest full
prefix pole, and use a finite coefficient minor of those
residues with the disjoint tail variables.  Lower initial
token weights cannot have this full prefix pole.
Descending through its Jordan orders removes that token.
Induction on the tail degree and then the initial weight
finishes.  The initial \(W\)-term uses ordinary completed
tensor injection.  This proves unrestricted injectivity
in every degree while establishing the induction step.

It suffices first to prove the quotient map modulo
\(D_\beta\) is injective in the range \(d<n\leq2d\).

\paragraph{Interior new tokens.}
A new weight-\(d\) token in an interior interval has
nonempty outside blocks of lengths \(a,b\), with
\(a+b=n-d\).  Its full interval has disconnected complement,
so its pole cannot coincide with a distinct interval pole
on the total hypersurface.  The outside source module is
the actual \(\mathcal A_a\widehat\otimes\mathcal A_b\).
Both \(a,b<d\), so
Lemma~\ref{lem:wcs-balanced} injects its quotient at outside
total \(\beta/\alpha\).  The independent internal companion
residues therefore force its outside coefficient to be an
actual source multiple of
\(\gamma_{\rm out}-\beta/\alpha\).

For a Jordan chain \(\gamma y=\alpha y+Ny\), the source
coinvariant relation is exactly
\begin{equation}\label{eq:wcs-Jordan}
 N^k y\otimes(\gamma_{\rm out}-\beta/\alpha)b
 \equiv-\alpha^{-1}N^{k+1}y\otimes\gamma_{\rm out}b
                                      \pmod{\gamma-\beta}.
\end{equation}
At the end of the chain its companion belongs to the old
algebra.  Repeating removes the interior placement by
actual source relations.  The process is finite in the
number of positions and the Jordan lengths.  In particular
it remains valid when \(n=2d\) and \(a+b=d\).

\paragraph{A boundary token when \(d<n<2d\).}
Put \(s=n-d\).  Modulo old prefix terms, the contribution
beginning in \(V_d\) has the actual balanced form
\begin{equation}\label{eq:wcs-new-prefix}
 \sum_i f^{-r_i}C_i(\text{prefix},f)
                       \otimes b_i(\text{suffix},f),
 \qquad b_i\in E_s .
\end{equation}
Use the local suffix decomposition
\[
 b_i=b_i^\circ+\sum_a f^{-q_a}D_a\,b_{ia},
 \qquad b_i^\circ\in S_s^\circ,\quad b_{ia}\in R_f.
\]
Here \(S_s\) is the algebra in degree \(s\) \emph{before}
weight \(s\) was adjoined.  Its total saturation is part of
the retained smaller-alphabet induction, even though
\(s\) is below the current old maximum weight.

Apply the suffix companion-column projectors first.
By Lemma~\ref{lem:wcs-boundary-bound}, every other term
has zero negative principal part in the complementary
coordinates.  A finite prefix minor of the independent
\(C_i\) columns therefore shows
\[
 \iota(b_i^\circ)\in f^{r_i}\mathcal B_s.
 \quad\text{Hence}\quad b_i^\circ\in f^{r_i}S_s^\circ
\]
by old saturation.  On the \(a\)-th polar line, every
other term has order at most \(q_a\), whereas
\eqref{eq:wcs-new-prefix} gives
\(\sum_i f^{-r_i-q_a}C_i b_{ia}\).  Multiply by \(f^{q_a}\)
and use the same prefix minor.  It gives
\(b_{ia}\in f^{r_i}R_f\).  Thus the exact source conclusion is
\begin{equation}\label{eq:wcs-tail-divisibility}
                         b_i\in f^{r_i}E_s .
\end{equation}

Balancing sends \(f=r-\alpha\) to
\(\beta\gamma_s^{-1}-\alpha\), a unit multiple of
\(\gamma_s-\beta/\alpha\).  On this Jordan block the
quotient of the new prefix by the old prefix is exactly
\[
 E_s/(\gamma_s-\beta/\alpha)^{r_i}E_s .
 \]
Equation~\eqref{eq:wcs-tail-divisibility} is zero in this
source quotient.  Its finitely many root-local conclusions
give global zero by Bezout, since the quotient is killed by
the finite transported characteristic polynomial of \(V_d\).
Right exactness of the actual coinvariant sequence now
replaces the prefix contribution by an old source element
plus an actual \((\gamma-\beta)\)-multiple.  The companion
correction has no weight-\(d\) token.

The symmetric argument treats the suffix.  At most one
new token occurs for \(n<2d\).  Having removed its interior
and boundary placements, the remaining source is old;
old degree-\(n\) saturation finishes this range.

\paragraph{The paired boundary case \(n=2d\).}
After removing interior placements, all remaining new
terms lie in the single cut module
\[
                  (E_d\widehat\otimes E_d)_{\gamma=\beta}.
\]
This includes both edge orientations and the possible
two-token word \(V_dV_d\).  Every term outside this cut
has no full prefix-\(d\) or suffix-\(d\) pole.

Use the finite polar coordinates in both factors, with
the second total transported by \(t=\beta/r\).
The old balanced module
\((S_d\widehat\otimes S_d)_{\gamma=\beta}\)
is saturated at \(f=r-\alpha\): modulo \(f\) its map is
the completed tensor of the two old specialized
injections in disjoint internal variables.  Repeated
division gives every power.
The new tensor coordinates are rescaled by
\(f^{-r_i-q_a}\), with exponent zero on each unchanged
line.  Since the total cut coefficient is regular,
its coordinate numerators are divisible by the full
rescaling powers.  It belongs to the old balanced
source locally at every relevant root.  The quotient
is killed by the product of the two companion
polynomials, one transported by \(r\mapsto\beta/r\);
Bezout gives global old-cut membership.  This is an
actual source replacement.  Old saturation finishes
degree \(2d\).

\paragraph{Degrees \(n>2d\).}
Filter by the largest permitted weight \(j\) of the
first token.  Its full prefix pole becomes a total
divisor of the tail of length \(n-j>d\).  No proper
companion interval has that length, and a lower
initial token cannot contain the full prefix.
Thus the highest Jordan residue is separate from the
lower filtration.  Finite coefficient extraction in
the prefix and the already proved lower-degree total
saturation put its tail coefficient in the actual
source image of \(\gamma_{\rm tail}-\beta/\alpha\).
Equation~\eqref{eq:wcs-Jordan} removes it.  This descends
through the first-token weights.

At the regular first token, the exact source evaluation is
\[
 z_1=\beta\gamma_{\rm tail}^{-1}.
\]
For \(X_1=z_1-1\), it sends a series
\(\sum_m X_1^m b_m\) to
\(\sum_m(\beta\gamma_{\rm tail}^{-1}-1)^m b_m\).
The operator in parentheses is topologically nilpotent
on every finite integral quotient, so the series
converges there.  Its kernel is the image of
\(X_1-(\beta\gamma_{\rm tail}^{-1}-1)\), by the exact
identity \(X^m-A^m=(X-A)\sum_{\ell<m}X^{m-1-\ell}A^\ell\)
and continuity.  Unrestricted tail injection finishes
this step.  This proves all \(n>d\).

\paragraph{Powers and primitive elements.}
If \(\iota_n(a)\) is divisible by \(D_\beta^k\), the
first-power assertion gives \(a=(\gamma-\beta)a_1\)
in the actual source.  Since multiplication by
\(D_\beta\) is injective in the coefficient domain,
\(\iota_n(a_1)\) is divisible by \(D_\beta^{k-1}\).
Repeating proves \eqref{eq:wcs-saturation}.  Each fixed
\(n,k\) uses finitely many bounded divisions; the final
source element has a common denominator.

For primitive companions, the Eulerian idempotent
\[
 e=\log^\star(\operatorname{id})
 \]
projects the completed tensor algebra onto its primitive
free Lie algebra; see \cite[\S1.1, p.~2]{BandieraSchaetz}.
In fixed weighted degree this is a finite sum with a
common rational denominator.  Its naturality makes it
commute with the Hopf automorphism \(\gamma\), even when
\(\gamma\) changes ordinary token count.  Apply \(e\)
to the source equality
\(\ell=(\gamma-\beta)^k a\) for primitive \(\ell\).
It gives
\(\ell=(\gamma-\beta)^k e(a)\) in the free weighted
Lie source, which proves the last assertion.
\end{proof}

\subsection{The exact cut relation represented by the proof}

For \(n>d\), the actual source maps
\[
 \bigoplus_{i+j+k=n}\mathcal A_i\widehat\otimes
       \mathcal A_j\widehat\otimes\mathcal A_k
 \xrightarrow{\ \partial\ }
 \bigoplus_{i+j=n}\mathcal A_i\widehat\otimes\mathcal A_j
 \longrightarrow\mathcal A_n\longrightarrow0,
 \qquad
 \partial(a\otimes b\otimes c)=ab\otimes c-a\otimes bc,
\]
are exact.  For each word, its possible cuts form a
nonempty simplex, and these are exactly its cut-identification
relations.  The coefficientwise contraction has coefficients
\(0,1,-1\), so it is bounded on the completions.
Every term and every cycle is \((\gamma-\beta)\)-torsion-free
by unrestricted injection.  The exact sequences of cycles
therefore remain exact on coinvariants.

The proof above shows that every relation among the
specialized coefficient images of all cuts is the image
of this actual three-block boundary.  The boundary
principal-part argument is needed here: intersections
of two specialized cut images can be mediated by a
third cut, and a crossing token can carry a complementary
internal pole.  Equations~\eqref{eq:wcs-ancestor} and
\eqref{eq:wcs-tail-divisibility} give their source
replacement without discarding either phenomenon.
